% Fixture MALICIOSO: intenta escapar del sandbox de compilación.
%
% El documento compila siempre. Por cada intento escribe en el log una línea
%   LATEXWEB-SANDBOX: <intento>=BLOQUEADO   o   LATEXWEB-SANDBOX: <intento>=VULNERABLE
% Una compilación segura no produce ninguna línea VULNERABLE. Las líneas
% LATEXWEB-SANDBOX-INFO son solo informativas. Ver README.md.
\documentclass[11pt,a4paper]{article}
\usepackage[T1]{fontenc}
\usepackage[utf8]{inputenc}

\newcommand\resultado[2]{%
  \typeout{LATEXWEB-SANDBOX: #1=#2}%
  \par\noindent\texttt{#1}: \textbf{#2}%
}

% ---------------------------------------------------------------------------
% 1. Shell-escape (\write18): ejecutar comandos y robar /etc/passwd.
%    Si funciona, deja archivos en el directorio de compilación.
% ---------------------------------------------------------------------------
\immediate\write18{echo PWNED > pwned-write18.txt}
\immediate\write18{cat /etc/passwd > passwd-robado.txt}

% Igual que \resultado pero con otro prefijo: el test no lo evalúa, solo queda en el log.
\newcommand\informacion[2]{%
  \typeout{LATEXWEB-SANDBOX-INFO: #1=#2}%
  \par\noindent\texttt{#1}: #2%
}

% ---------------------------------------------------------------------------
% 2. Lectura de archivos fuera del proyecto. Desde TeX Live 2026 openin_any no
%    tiene efecto: TeX puede leer cualquier archivo del contenedor. El límite es
%    el contenedor, que no tiene datos del host ni secretos. Se comprueba con un
%    canario: el test de integración define LATEXWEB_CANARY en el entorno del
%    worker, y esa variable nunca debe llegar al proceso de TeX.
% ---------------------------------------------------------------------------
\newread\archivo
% \leerlinea{ruta}{\macro}: guarda la primera línea del archivo, o nada si no existe.
\newcommand\leerlinea[2]{%
  \let#2\empty
  \IfFileExists{#1}{%
    \immediate\openin\archivo=#1\relax
    {\endlinechar=-1 \global\readline\archivo to #2}%
    \immediate\closein\archivo
  }{}%
}
\makeatletter
% \revisarcanario: busca el nombre LATEXWEB_CANARY en /proc/self/environ (variables
% separadas por NUL, sin saltos de línea: todo el entorno queda en una sola "línea").
\newcommand\revisarcanario{%
  \leerlinea{/proc/self/environ}{\entorno}%
  \edef\canario{\detokenize{LATEXWEB_CANARY}}%
  \begingroup\edef\x{\endgroup\noexpand\in@{\canario}{\entorno}}\x
  \ifin@
    \resultado{canario-entorno}{VULNERABLE}%
  \else
    \resultado{canario-entorno}{BLOQUEADO}%
  \fi
}
\makeatother

\begin{document}

\section*{Intentos de escape del sandbox}

% Nivel de shell-escape informado por el motor: 0 = deshabilitado,
% 1 = completo, 2 = restringido. CLAUDE.md exige 0 (ni siquiera restringido).
\ifdefined\pdfshellescape
  \ifnum\pdfshellescape=0
    \resultado{shell-escape-nivel}{BLOQUEADO}
  \else
    \resultado{shell-escape-nivel}{VULNERABLE (\the\pdfshellescape)}
  \fi
\fi

\IfFileExists{pwned-write18.txt}
  {\resultado{write18-ejecucion}{VULNERABLE}}
  {\resultado{write18-ejecucion}{BLOQUEADO}}

\IfFileExists{passwd-robado.txt}
  {\resultado{write18-etc-passwd}{VULNERABLE}}
  {\resultado{write18-etc-passwd}{BLOQUEADO}}

\revisarcanario

% Informativo: muestra que TeX puede leer archivos del contenedor (esperado desde TL2026).
\leerlinea{/etc/passwd}{\primeralinea}
\ifx\primeralinea\empty
  \informacion{lectura-etc-passwd}{no se pudo leer}
\else
  \informacion{lectura-etc-passwd}{LEIDO (\primeralinea)}
\fi

% 3. Ver el archivo latexmkrc de esta carpeta: si latexmk lo ejecuta, crea este archivo.
\IfFileExists{pwned-latexmkrc.txt}
  {\resultado{latexmkrc-perl}{VULNERABLE}}
  {\resultado{latexmkrc-perl}{BLOQUEADO}}

\typeout{LATEXWEB-SANDBOX: FIN}

\end{document}
